\section{The asymptotically optimal coefficient for few facets}\label{sec:few}
\subsection{From the actual facet law to a single clipping}
Proposition~\ref{prop:enclosure}, proved in full in the appendices, gives
the precise finite-facet interface needed here: for
$0\le e(K)\le H_d^{-1}$ there is an enclosing simplex $P$, formed by
exactly $d+1$ distinct irredundant $K$-facet halfspaces, such that
\begin{equation}\label{eq:facet-proj}
 \pi_P(u)-\pi_K(u)\le H_de(K)\pi_P(u)\quad(\norm u=1).
\end{equation}
It uses the actual atoms $X_F=\nu_F/h_F$ and probabilities
$p_F=s_Fh_F/(d|K|)$. In particular, the selected halfspaces are actual
facets, rather than unspecified support planes inferred from a compact
limit. The determinant identity is exactly
$(B-A)/B=ne/(1+ne)$ for this law. The assignment bound is
$h\le n^3(n+1)e/(1+ne)$; at the upper endpoint
$h\le1/(3d^2)<1/(2d)$, so its strict enclosure gate remains valid.

We also record the projection-cap argument explicitly.
\begin{lemma}[Near corners from all projections]\label{lem:cap}
If $K\subset P$, with $P$ a full-dimensional simplex, and
$\pi_P(u)-\pi_K(u)\le\delta\pi_P(u)$ for every unit $u$, where
$0\le\delta\le1$, then each corner loss satisfies
\[
 1-\max_{x\in K}\lambda_i^P(x)\le\delta^{1/(d-1)}.
\]
\end{lemma}
\begin{proof}
Let $q_i$ be the corresponding vertex and let
$t_i=1-\max_K\lambda_i^P\in[0,1]$. Choose a unit direction parallel to
its opposite facet. The affine function $\lambda_i^P$ descends under
this projection, and has range $[0,1]$ on the projected simplex.
For $0<s<t_i$, the projected homothetic cap $q_i+s(P-q_i)$ is contained
in $\{\lambda_i^P\ge1-s\}$ and disjoint from the projection of $K$.
Its projected volume is $s^{d-1}\pi_P(u)$. Both sets lie in the
projection of $P$, so $s^{d-1}\le\delta$. Letting $s\uparrow t_i$
proves the claim; $t_i=0$ is immediate. Strict caps avoid any assumption
about the measure of the boundary, and include $\delta=0$.
\end{proof}

\begin{proof}[Proof of the upper bound in Theorem~\ref{thm:few-facet}]
A full-dimensional polytope has at least $d+1$ facets. With exactly $d+1$
it is a simplex, and $E=e=0$. Otherwise it has $d+2$ facets. The $P$
above uses $d+1$ of them, so $K=P\cap H$ for the remaining facet
halfspace. Set $\rho=(H_de(K))^{1/m}$. Lemma~\ref{lem:cap} bounds all
corner losses by $\rho$, and Theorem~\ref{thm:clipping} gives
$E(K,S)\le n\rho$ for this same originally prescribed maximum.
The argument includes $e=0$ and $e=H_d^{-1}$. For larger $e$, the universal
maximum-simplex bound $E\le n$ suffices, since $\rho>1$.

For the uniform scalar estimate, $H_3=8640\le144^2$, and
$H_{d+1}/H_d\le18\le64$ follows from
$(d+1)/d\le2$, $(d+2)/(d+1)\le3/2$, and
$(d+3)/(d+2)\le4/3$ in its squared, cubed, and linear factors.
Induction gives $H_d\le144^{d-1}$. Since $n\le4d/3$,
$nH_d^{1/m}\le192d$.
\end{proof}

\subsection{Sharp exponent and a realized lower coefficient}
The body $K_t$ in \eqref{eq:truncation-body} has exactly $d+2$ facets.
Theorem~\ref{thm:clipping} identifies vertex maxima with
$E(K_t,S_t)=nt$, and \eqref{eq:truncation-exact} gives
$e(K_t)\sim d(d-1)t^m/n^2$. Thus for every fixed $d$ and every
$\eta>1/m$, the ratio $E/e^\eta$ diverges as $t\downarrow0$.
No such stronger exponent can hold uniformly even with a
dimension-dependent coefficient.

For the all-scale lower coefficient take
\[
 C_d=\conv(0,e_1,e_2,e_1+e_2,e_3,\ldots,e_d),\qquad
 S_d=\conv(0,e_1,\ldots,e_d).
\]
This is a square followed by $d-2$ pyramids, with irredundant inequalities
\[
 x_i\ge0\ (1\le i\le d),\qquad
 x_1+\sum_{i=3}^dx_i\le1,\qquad x_2+\sum_{i=3}^dx_i\le1.
\]
Every full-dimensional vertex simplex contains all pyramid vertices and
three of the four square vertices. Every such square triangle has area
$1/2$, so the resulting volume is $1/d!$. Separate determinant expansion
in convex combinations bounds arbitrary inscribed simplices by the same
volume. Thus $S_d$ is maximum. The omitted square vertex has coordinates
$(-1,1,1,0,\ldots,0)$ in $S_d$, so $E(C_d,S_d)=n$.

For completeness, a direct actual-facet certificate gives the defect.
Set $c=1/(d-1)!$. The area-normal/support pairs are
\[
 (-ce_1,0),\quad(-ce_2,0),\quad(-2ce_i,0)\ (3\le i\le d),
\]
\[
 \left(c\left(e_1+\sum_{i=3}^de_i\right),c\right),\qquad
 \left(c\left(e_2+\sum_{i=3}^de_i\right),c\right).
\]
The volume is $2/d!$. For the raw sums of
Appendix~\ref{prop:actual-moments}, direct minors give
\[
 \mathcal H(C_d)=c^d d2^{d-1},\qquad
 \mathcal L(C_d)=c^{d+1}2^d.
\]
Only four lifted subsets contribute, each by $c^{d+1}2^{d-2}$:
omit one top facet or one of the first two coordinate facets.
For horizontal minors the contributions with zero, one, and two top
facets are respectively $1,d,d-1$, multiplied by $c^d2^{d-2}$.
The exact actual identity $a=\mathcal L/(d|C_d|\mathcal H)$ therefore
gives $a(C_d)=1/d$ and $e(C_d)=1/(dn)$. Hence
$G_d^{\rm fac}\ge n[dn]^{1/m}$.

Finally, $H_d$ is a fixed polynomial in $d$, so $H_d^{1/m}\to1$;
also $[dn]^{1/m}\to1$. Divide the two coefficient bounds by $d$ and
squeeze to obtain \eqref{eq:few-squeeze}. No uniformity in a
fixed-dimensional truncation remainder is needed.
